\chapter{Turbulence, diffusion and damping}\label{ch:turb}

Motions smaller than the grid spacing transport momentum, heat and moisture. WRF represents them by eddy
diffusivities that act on the resolved fields. Which closure applies depends on the resolution:
\begin{itemize}
  \item on d01 (900~m), the boundary layer is parameterized by YSU (\cref{ch:pbl}), and horizontal mixing uses a
    two-dimensional Smagorinsky closure;
  \item on d02 (100~m), a three-dimensional large-eddy closure with a prognostic turbulent kinetic energy (TKE)
    resolves the large eddies of the boundary layer.
\end{itemize}
This chapter describes both closures, the stress and flux computations they share, and the damping layers that keep
the model stable. The routines are in \code{dyn_em/module_diffusion_em.F}; \code{first_rk_step_part2} calls them in
RK stage 1.

\section{Options}

\begin{description}
  \item[\code{diff_opt}] How diffusion is computed. Option 2, used here, evaluates the full stress tensor in physical
    space, with the terrain metric terms. Option 1 diffuses along coordinate surfaces.
  \item[\code{km_opt}] How the eddy coefficients are computed:
    \begin{itemize}
      \item 2: a 1.5-order closure with prognostic TKE (d02);
      \item 4: horizontal two-dimensional Smagorinsky, vertical mixing from the PBL scheme (d01);
      \item 3: three-dimensional Smagorinsky; 5: a scale-aware variant.
    \end{itemize}
\end{description}

\section{The deformation tensor}

Both closures use the deformation (strain-rate) tensor $D_{ij} = \partial_j u_i + \partial_i u_j$. Its components are
computed on their natural staggered positions by \src{dyn_em/module_diffusion_em.F}{cal_deform_and_div}:
\begin{itemize}
  \item the diagonal $D_{11}$, $D_{22}$, $D_{33}$ at mass points;
  \item $D_{12}$ at vorticity points;
  \item $D_{13}$, $D_{23}$ at $u$-$w$ and $v$-$w$ edges.
\end{itemize}
With terrain, horizontal derivatives on a coordinate surface are corrected by the slope of the surface:
\begin{equation}
  \left.\pd{\psi}{x}\right|_z = \left.\pd{\psi}{x}\right|_\eta - \frac{z_x}{z_\eta}\,\pd{\psi}{\eta},
\end{equation}
where $z_x$ and $z_\eta$ are the metric terms computed once per stage by
\src{dyn_em/module_diffusion_em.F}{compute_diff_metrics}. That routine also computes the heights of the levels and
their inverse spacings.

\section{Two-dimensional Smagorinsky (d01)}

The horizontal eddy viscosity is
\begin{equation}
  K_h = (C_s\,l)^2\sqrt{\tfrac14(D_{11} - D_{22})^2 + D_{12}^2}, \qquad l = \sqrt{\Delta x\,\Delta y},
\end{equation}
with $C_s = 0.25$ (\code{c_s}), in \src{dyn_em/module_diffusion_em.F}{smag2d_km}. Over steep terrain $K_h$ is reduced
by the local slope factor $\alpha = \max(\sqrt{z_x^2 + z_y^2}, 1)$: divided by $\alpha^2$ where the deformation is
large and by $\alpha$ elsewhere. The heat diffusivity is $K_h/\mathrm{Pr}$ with $\mathrm{Pr} = 1/3$
(\code{prandtl}). Vertical mixing is left to the PBL scheme.

\section{The 1.5-order TKE closure (d02)}

The subgrid kinetic energy $e$ has its own prognostic equation:
\begin{equation}
  \partial_t(\mu_d e) + \text{advection} = \mu_d\,(P_s + P_b - \epsilon) + \text{diffusion},
\end{equation}
with shear production $P_s$, buoyancy production $P_b$ and dissipation $\epsilon$:
\begin{equation}
  P_s = K_h(D_{11}^2 + D_{22}^2 + D_{12}^2) + K_v(D_{33}^2 + D_{13}^2 + D_{23}^2) \;(\text{schematically}),\qquad
  P_b = -K_{hv}\,N^2, \qquad
  \epsilon = C\,\frac{e^{3/2}}{l}.
\end{equation}
The routines \src{dyn_em/module_diffusion_em.F}{tke_shear}, \code{tke_buoyancy} and \code{tke_dissip} evaluate them.
\code{calculate_N2} computes the Brunt--V\"ais\"al\"a frequency squared, including the moist form in saturated
air.

\paragraph{Length scales and coefficients.} In the anisotropic form used here (\code{mix_isotropic = 0}),
\src{dyn_em/module_diffusion_em.F}{tke_km} sets:
\begin{align}
  l_h &= \sqrt{\Delta x\,\Delta y}, \qquad
  l_v = \begin{cases} \min\bigl(\Delta z,\; 0.76\,\sqrt{e}/N\bigr) & N^2 > 0,\\ \Delta z & \text{otherwise},\end{cases}\\
  K_{mh} &= c_k\,\sqrt{e}\;l_h, \qquad K_{mv} = c_k\,\sqrt{e}\;l_v, \qquad c_k = 0.15,\\
  K_{hh} &= K_{mh}/\mathrm{Pr}, \qquad K_{hv} = K_{mv}\,\bigl(1 + 2\,l_v/\Delta z\bigr),
\end{align}
with a lower bound proportional to $l^2$ and an upper bound \code{mix_upper_bound}$\cdot l^2/\Delta t$ for
stability. For the dissipation, \code{calc_l_scale} uses the isotropic scale
$\Delta s = (\Delta x\,\Delta y\,\Delta z)^{1/3}$, reduced in stable air:
\begin{equation}
  l = \min\Bigl(\Delta s,\ \max\bigl(0.76\,\sqrt{e}/N,\ 0.001\,\Delta s\bigr)\Bigr) \quad (N^2 > 10^{-6}),
\end{equation}
and the coefficient of $\epsilon$ is
\begin{equation}
  C = c_{e1} + c_{e2}\,\frac{l}{\Delta s}, \qquad c_{e1} = \frac{c_k}{0.10}\,0.19, \qquad c_{e2} = \max(0, 0.93 - c_{e1}),
\end{equation}
with $C = 3.9$ at the lowest and highest levels. The cube root and the power $e^{3/2}$ are evaluated with the
reproducible \code{rp_pow} (\cref{ch:repro}).

\section{Fluxes and tendencies}

With the eddy coefficients, the stresses are
\begin{equation}
  \tau_{ij} = -\mu_d\,K\,D_{ij},
\end{equation}
computed by \code{cal_titau_11_22_33}, \code{cal_titau_12_21}, \code{cal_titau_13_31} and \code{cal_titau_23_32}.
Their divergence gives the momentum tendencies:
\begin{itemize}
  \item horizontal: \code{horizontal_diffusion_u_2}, \code{_v_2}, \code{_w_2};
  \item vertical: \code{vertical_diffusion_u_2}, \code{_v_2}, \code{_w_2}.
\end{itemize}
Scalars (θ, moisture, TKE) diffuse with the heat diffusivities: \code{horizontal_diffusion_s} and
\code{vertical_diffusion_s}.
\begin{itemize}
  \item At the surface, the vertical fluxes are not diffusive: they come from the surface layer scheme
    (\cref{ch:sfc}), as stresses $u_*^2$ and fluxes of heat and moisture.
  \item With \code{m_opt = 1}, the stresses are also kept for output (\code{nba_mij}).
\end{itemize}

\begin{portnote}
\code{cal_deform_and_div} and the horizontal diffusion routines are long sequences of short loop nests over the whole
domain: dozens of nests, each reading and writing whole 3D fields. On the GPU each nest is one kernel, in source
order (59 launches for \code{cal_deform_and_div} alone in the port, \cref{ch:subsys}), so this part of the step is
dominated by memory traffic and launches. It is a prime candidate for fusion and tiling (\cref{ch:fusion,ch:tiling}).
\end{portnote}

\section{Damping}

\paragraph{Rayleigh damping of $w$ at the top} (\code{damp_opt = 3}). In the top \code{zdamp} metres (5000~m by
default), $w$ is relaxed toward zero inside the implicit acoustic solve, in \code{advance_w}. The weight is
\begin{equation}
  d(z) = \Delta\tau\,\gamma_r\,\sin^2\!\left(\frac{\pi}{2}\,\frac{z - z_b}{z_{\mathrm{damp}}}\right) \quad (z > z_b),
  \qquad W \gets \frac{W - d\,\mu_f\,W^{\mathrm{old}}}{1 + d},
\end{equation}
where $z_b$ is the bottom of the layer and $\gamma_r$ is \code{dampcoef} (0.2~s$^{-1}$). The square of the sine is
written as two calls of \code{rp_sin} multiplied, as in the source. Damping absorbs gravity waves that would reflect
from the rigid lid.

\paragraph{Vertical-velocity damping} (\code{w_damping = 1}). Where the vertical Courant number exceeds a threshold
(\code{w_crit_cfl}), \src{dyn_em/module_big_step_utilities_em.F}{w_damp} adds a damping tendency to $W$. This
prevents instability in strong updrafts, such as those over an intense fire. The routine also reports the maximum
Courant numbers; the port computes them with exact maximum reductions.

\paragraph{Divergence damping and the external-mode filter} act in the acoustic steps (\cref{ch:time}).
